

\section{La fausse solution des LCA (2016) : Dilution de la langue dans la "culture"}
space{0.3cm}
oindent

L'histoire de l'enseignement des Langues et Cultures de l'Antiquité (LCA) au sein du système
éducatif français est marquée par une succession de crises de légitimité et de réformes
structurelles, mais la réforme du collège de 2016, portée par la ministre Najat Vallaud-Belkacem
dans le cadre de la « Refondation de l'École », constitue une rupture épistémologique et
institutionnelle sans précédent. Souvent qualifiée de « cataclysme » par les praticiens et les
associations de spécialistes telles que la Coordination Nationale des Associations Régionales des
Enseignants de Langues Anciennes (CNARELA) ou le Syndicat National des Enseignements de Second degré
(SNES-FSU) 1, cette réforme ne s'est pas contentée d'ajuster des volumes horaires ; elle a redéfini
la nature même de l'objet enseigné.

L'objet de ce rapport de recherche est d'analyser en profondeur ce que l'on peut qualifier de «
fausse solution » de 2016. Cette expression désigne le compromis politique et pédagogique bancal qui
a prétendu « sauver » les langues anciennes en les diluant dans une approche culturaliste
interdisciplinaire — les Enseignements Pratiques Interdisciplinaires (EPI) — tout en réduisant les
horaires disciplinaires à des volumes « homéopathiques » qui rendent l'apprentissage linguistique
inopérant pour le plus grand nombre. Cette configuration a engendré un double échec structurel.
D'une part, la « dilution culturelle » a affaibli la spécificité de la discipline, la transformant
en une variable d'ajustement budgétaire et pédagogique au sein des établissements. D'autre part, la
réduction drastique des heures a paradoxalement figé les pratiques pédagogiques : faute de temps
pour l'immersion, l'automatisation ou les méthodes actives prônées par la recherche didactique
récente (comme le projet ELLASS), la méthode traditionnelle « grammaire-traduction » a perduré par
défaut. Cette rémanence méthodologique, appliquée sur des horaires squelettiques, a créé une
surcharge cognitive insupportable pour les élèves, conduisant à un élitisme de fait, contraire aux
objectifs affichés de démocratisation de la réforme.

Ce document, exhaustif et rigoureusement documenté, s'appuie sur une analyse critique des textes
officiels (Bulletin Officiel de l'Éducation Nationale, arrêtés, décrets), des rapports d'inspection
générale (IGESR), de la littérature didactique contemporaine (travaux sur la charge cognitive
d'André Tricot et John Sweller), ainsi que des prises de position syndicales et associatives. Il
explorera successivement la déstructuration institutionnelle opérée en 2016, la mutation
épistémologique vers une « culture » sans langue, l'impasse didactique de la grammaire-traduction
dans des horaires contraints, et enfin les conséquences à long terme de cette réforme, jusqu'au
récent « Choc des savoirs » de 2024 qui semble porter le coup de grâce à cet édifice fragilisé.

Pour saisir la portée de la « fausse solution », il est impératif de disséquer l'architecture
administrative qui l'a rendue possible. La réforme de 2016 n'a pas seulement réduit les moyens ;
elle a opéré une mutation statutaire des Langues et Cultures de l'Antiquité, les faisant basculer
d'une discipline optionnelle fléchée à un dispositif périphérique et précaire.



\subsection{De la discipline canonique à l'enseignement de complément : une dégradation statutaire systémique}

Avant la réforme de 2016, l'enseignement du latin et du grec bénéficiait, malgré des érosions
successives au fil des décennies, d'un statut d'option relativement protégé et identifié dans la
grille horaire nationale. Les élèves choisissaient l'option, et les heures étaient dues par
l'institution. La réforme du collège de 2016 a introduit une rupture sémantique et juridique majeure
en redéfinissant les LCA non plus comme une discipline à part entière dans la grille horaire
standard, mais comme un « Enseignement de Complément » adossé à des Enseignements Pratiques
Interdisciplinaires (EPI).\cite{I-3-2-ref1}



\subsection{La mécanique des horaires « homéopathiques » et l'effondrement du temps d'exposition}

Le terme « homéopathique » qualifie avec une justesse cruelle les volumes horaires entérinés par
l'arrêté du 19 mai 2015 relatif à l'organisation des enseignements au collège. Alors que les
horaires traditionnels permettaient une exposition régulière (souvent 2 heures en 5ème et 3 heures
dès la 4ème), la nouvelle grille a imposé un rationnement drastique qui place l'enseignement en deçà
du seuil d'efficacité pédagogique.

L'analyse comparée des textes officiels et de leur mise en application révèle la structure suivante
pour le cycle 4 (5ème, 4ème, 3ème) :

\begin{itemize}\item \textbf{Classe de 5ème :} La dotation a été réduite à 1 heure hebdomadaire.\cite{I-3-2-ref1} Ce volume est considéré par la quasi-totalité des didacticiens et des praticiens comme inopérant pour l'apprentissage d'une langue flexionnelle. Une heure par semaine ne permet ni l'installation de rituels d'apprentissage, ni la mémorisation efficace (la courbe de l'oubli opérant pleinement entre deux séances espacées de sept jours), ni l'immersion culturelle. Il s'agit d'une logique d'affichage politique : le ministère peut clamer qu'il y a « du latin dès la 5ème », mais la réalité didactique est celle d'une initiation superficielle.\item \textbf{Classe de 4ème :} L'horaire officiel est passé à 2 heures hebdomadaires, avec une possibilité théorique d'extension à 3 heures, mais cette extension est soumise à des contraintes locales fortes.\cite{I-3-2-ref1}\item \textbf{Classe de 3ème :} De même, l'horaire est fixé à 2 heures, extensible à 3 heures.\cite{I-3-2-ref1}\end{itemize}Ce volume global, oscillant entre 1 et 2 heures par semaine, constitue une régression historique.
Comme le soulignent les syndicats enseignants, cet horaire réduit a été vécu comme une tentative de
« suppression pure et simple » déguisée en maintien symbolique.\cite{I-3-2-ref1} L'association
\textit{Arrête Ton Char} a démontré dans ses mémorandums que cette réduction horaire rendait
impossible l'atteinte des objectifs linguistiques fixés par ailleurs dans les programmes, créant une
injonction paradoxale permanente pour les enseignants.\cite{I-3-2-ref5}



\subsection{Le piège sémantique et juridique : « dans la limite de »}

Une subtilité administrative majeure, lourde de conséquences, réside dans la formulation des textes
réglementaires ajustés après la fronde de 2015-2016. L'arrêté du 16 juin 2017, venu modifier celui
de 2015, précise que les enseignements facultatifs peuvent porter sur les LCA « dans la limite d'une
heure hebdomadaire en classe de cinquième et de trois heures hebdomadaires pour les classes de
quatrième et de troisième ».\cite{I-3-2-ref3}

Cette locution « dans la limite de » est cruciale. Elle opère un renversement de la logique de droit
: l'horaire n'est plus un plancher garanti (un droit de l'élève), mais un plafond autorisé (une
variable de gestion).

\begin{itemize}\item \textbf{Plafond vs Plancher :} En droit administratif scolaire, ce qui est défini comme une limite supérieure devient rarement la norme effective, surtout en période de contrainte budgétaire.\item \textbf{Discrétionnarité locale :} Ce libellé a transféré la responsabilité de l'horaire du niveau national (le Ministère) au niveau local (le Conseil d'Administration du collège). Ce sont les chefs d'établissement, contraints par des Dotations Horaires Globales (DHG) insuffisantes, qui doivent décider d'attribuer 1h, 2h ou 3h.\item \textbf{Variable d'ajustement :} En pratique, cela a permis aux administrations locales de ne proposer que le minimum vital, voire de ne pas ouvrir l'option certaines années, en arguant de contraintes budgétaires ou de concurrence avec d'autres projets. Le SNES-FSU a relevé que cet horaire, théoriquement extensible à 3 heures, restait souvent bloqué à des niveaux inférieurs, l'extension nécessitant un financement sur la marge d'autonomie que peu d'établissements peuvent se permettre.\cite{I-3-2-ref1}\end{itemize}

\subsection{La guerre des moyens : la marge d'autonomie comme champ de bataille structurel}

L'un des aspects les plus pernicieux de la réforme de 2016 réside dans le mode de financement des
heures d'enseignement des LCA. Au lieu d'être sanctuarisées et financées par des heures postes
spécifiques attribuées par le rectorat en sus de la dotation de base, les heures de latin et de grec
ont été basculées sur la « marge d'autonomie » des établissements.\cite{I-3-2-ref1}



\subsection{Le mécanisme de concurrence interne et de culpabilisation}

La marge d'autonomie est une enveloppe d'heures (fixée à 2h45 puis 3h par division dans les textes
de 2016/2017) laissée à la discrétion du Conseil d'Administration du collège pour financer une
multitude de dispositifs pédagogiques concurrents :

1. Les groupes à effectifs réduits (dédoublements) en sciences, langues vivantes ou français.

2. La co-intervention (deux professeurs dans la même classe).

3. Les enseignements facultatifs (LCA, langues régionales, chorale, classes
bilangues).\cite{I-3-2-ref3}

Ce dispositif a placé structurellement les LCA en concurrence directe avec les besoins fondamentaux
des autres disciplines. Pour ouvrir un groupe de latin de 2 heures, un principal doit souvent
renoncer à dédoubler une classe de sciences expérimentales ou d'anglais surchargée. Comme l'analyse
la CNARELA, ce système crée une culpabilisation insidieuse des professeurs de Lettres Classiques :
demander l'ouverture de l'option ou le maintien d'un horaire décent revient à être accusé de «
prendre » les moyens des collègues pour une option jugée élitiste ou secondaire.\cite{I-3-2-ref1} La
solidarité des équipes pédagogiques est ainsi mise à l'épreuve par une pénurie organisée.



\subsection{L'iniquité territoriale et l'effet de taille}

L'association \textit{Arrête Ton Char} a produit des analyses quantitatives démontrant l'iniquité
intrinsèque de ce calcul de marge. La dotation est calculée par division (classe), mais les besoins
ne sont pas linéaires.

\begin{itemize}\item \textbf{L'effet de seuil :} Les petits collèges, ou ceux ayant des classes très chargées (au seuil de l'ouverture de classe), se retrouvent mathématiquement avec une marge de manœuvre plus faible pour financer des options coûteuses en heures.\item \textbf{Déserts de LCA :} Ce mécanisme favorise mécaniquement les gros établissements ou ceux situés dans des zones favorisées où la pression pour le maintien des options est forte et où les effectifs permettent de dégager de la marge. À l'inverse, dans les établissements les plus fragiles ou ruraux, le latin devient la première variable d'ajustement supprimée pour financer du soutien en français ou des dédoublements de sécurité, créant des déserts de LCA à rebours de l'objectif de démocratisation affiché par la réforme.\cite{I-3-2-ref5}\end{itemize}

\subsection{L'epi lca : une coquille vide et un leurre pédagogique?}

Pour compenser la perte d'heures disciplinaires et justifier la réforme, le ministère a mis en avant
la création des Enseignements Pratiques Interdisciplinaires (EPI), dont l'un des huit thèmes était «
Langues et Cultures de l'Antiquité ». Théoriquement, cet EPI devait permettre à tous les élèves du
cycle 4, et non plus seulement aux latinistes, d'accéder à la culture antique.\cite{I-3-2-ref1}

Cependant, l'analyse de la mise en œuvre sur le terrain révèle une « fausse solution » typique :

\begin{itemize}\item \textbf{Absence de cadrage national :} Le Conseil Supérieur des Programmes (CSP) n'a proposé aucun programme strict pour cet EPI, laissant les équipes locales improviser totalement le contenu.\cite{I-3-2-ref9}\item \textbf{Déspécialisation de l'enseignement :} L'EPI LCA pouvait légalement être enseigné par des professeurs d'autres disciplines (Histoire-Géographie, Français, Arts Plastiques, voire Sciences) sans aucune compétence linguistique avérée en latin ou grec ancien. Cela a dilué le contenu scientifique au profit d'une approche thématique souvent superficielle ou anecdotique.\cite{I-3-2-ref8}\item \textbf{Variable d'ajustement horaire :} L'EPI est souvent devenu un moyen administratif de justifier l'existence de l'option (« enseignement de complément ») sans lui donner les moyens d'un véritable cours de langue. Le SNES a souligné que l'EPI n'était pas un enseignement de latin, mais un projet vague, souvent ponctuel (sur un trimestre ou un semestre), incapable d'assurer la continuité nécessaire à l'apprentissage d'une langue.\cite{I-3-2-ref8}\end{itemize}

\subsection{Bilan quantitatif : une stabilisation en trompe-l'œil et une érosion réelle}

Les conséquences de cette ingénierie administrative ont été immédiates et durables. Même si le
ministère a parfois affiché des chiffres de stabilisation des effectifs globaux pour contrer les
critiques, la réalité du terrain rapportée par les syndicats et associations montre une
fragilisation profonde.

\begin{itemize}\item \textbf{Le trompe-l'œil statistique :} En 2017, une légère remontée des effectifs a été notée officiellement (passage de 401 498 à 415 987 latinistes).\cite{I-3-2-ref11} Cependant, cette hausse masque une baisse dramatique du volume horaire \textit{par élève}. Un élève qui suit 1 heure de latin par semaine est comptabilisé comme un « latiniste » au même titre que celui qui en suivait 3 heures auparavant. La \og couverture\fg{} s'est peut-être maintenue, mais l'intensité de la formation s'est effondrée.\item \textbf{L'effondrement du Grec Ancien :} Le grec ancien, discipline encore plus fragile, a subi une véritable hémorragie (1000 élèves en moins à la rentrée 2017).\cite{I-3-2-ref11} Son enseignement, nécessitant souvent des moyens spécifiques supplémentaires que la marge d'autonomie ne permettait plus de dégager face aux priorités du latin ou des autres matières, a été sacrifié dans de nombreux établissements.\item \textbf{Disparités d'accès :} L'enquête du SNES-FSU révèle que seulement 10\% des établissements ont bénéficié d'heures supplémentaires spécifiques pour les LCA, confirmant que pour la vaste majorité (90\%), l'enseignement s'est fait au détriment d'autres dispositifs ou a été réduit au minimum.\cite{I-3-2-ref12}\end{itemize}Au-delà de l'arithmétique des horaires, la réforme de 2016 a entériné un changement de paradigme
fondamental dans la conception même de la discipline. Le glissement sémantique officiel de « Latin »
ou « Grec » à « Langues et Cultures de l'Antiquité » (LCA) n'est pas anodin : il marque la priorité
donnée à la civilisation sur la langue, et à la compétence transversale sur le savoir disciplinaire
structuré.



\subsection{Analyse critique des programmes du cycle 4 (2016) : le triomphe du « culturel »}

Les programmes publiés au Bulletin Officiel en 2016 4 reflètent cette nouvelle philosophie
éducative. Ils sont structurés autour de compétences globales et de thématiques culturelles plutôt
que d'une progression grammaticale systématique et rigoureuse.



\subsection{La compétence culturelle comme alibi de la désubstantialisation linguistique}

Les textes officiels mettent l'accent de manière disproportionnée sur l'objectif « acquérir des
éléments de culture littéraire, historique et artistique ».\cite{I-3-2-ref13} Si cet objectif est
historiquement constitutif de la discipline, il tend ici à devenir prédominant, voire exclusif, face
à la réduction du temps imparti.

\begin{itemize}\item \textbf{Thématisation à outrance :} Les programmes sont découpés en grands thèmes (Mythes et héros, La République, Le Principat, Vie quotidienne, etc.) qui favorisent une approche documentaire.\cite{I-3-2-ref14} L'enseignant est incité à faire de la « civilisation » en français. Le texte latin, autrefois centre du cours, ne sert souvent plus que d'illustration, de vignette ou de prétexte à une activité de repérage lexical ponctuel.\item \textbf{Comparatisme et intertextualité :} L'accent est mis sur les liens entre l'Antiquité et le monde moderne, ou entre le français et le latin (étymologie, lexique savant).\cite{I-3-2-ref16} Cette approche, pertinente pour l'éveil culturel et la maîtrise du français, se substitue souvent à l'apprentissage rigoureux de la morphosyntaxe latine. On ne demande plus à l'élève de \textit{traduire} au sens strict de construction du sens par l'analyse, mais de « repérer et traiter les indices donnant accès au sens ».\cite{I-3-2-ref13}\end{itemize}Cette formulation officielle — « repérer des indices » — est révélatrice d'une démission didactique.
Elle valide une lecture de type devinette, une compréhension globale, intuitive et approximative, là
où la discipline exigeait traditionnellement une analyse précise des structures logiques. C'est ici
que réside la « dilution » : la langue n'est plus un système clos et cohérent à maîtriser, mais un
gisement d'indices culturels à exploiter superficiellement pour confirmer des connaissances acquises
par ailleurs (en cours d'histoire ou de français).



\subsection{La disparition de la traduction rigoureuse et de la progression linéaire}

Les programmes de 2016 encouragent une « progression en séquences décloisonnées ».\cite{I-3-2-ref9}
L'exercice canonique de la version ou du thème, jugé trop difficile ou élitiste, est relégué au
second plan au profit de la lecture cursive ou de la traduction accompagnée. L'objectif affiché est
de « développer des stratégies pour accéder au sens » sans nécessairement passer par l'analyse
exhaustive.\cite{I-3-2-ref13}

Cette évolution, motivée par le désir de ne pas rebuter les élèves et de rendre la matière «
attractive », a conduit selon les critiques (CNARELA, Association Guillaume Budé) à un « saupoudrage
» des connaissances. Sans maîtrise grammaticale, l'accès à la culture devient superficiel : l'élève
reste dépendant des traductions fournies par le professeur ou le manuel, en incapacité d'accéder
directement à la source textuelle. Il devient un consommateur de culture antique plutôt qu'un
lecteur de textes antiques.\cite{I-3-2-ref1}



\subsection{L'illusion de l'interdisciplinarité et la perte d'autonomie épistémologique}

L'EPI LCA est l'archétype de la dilution épistémologique. Conçu pour « construire une culture
commune », il a souvent abouti à une instrumentalisation du latin.

\begin{itemize}\item \textbf{Le latin au service des autres disciplines (ancillaire) :} Dans le cadre des EPI, le latin est souvent convoqué uniquement pour éclairer un point d'Histoire (la citoyenneté à Athènes) ou de Français (le vocabulaire du théâtre). Il perd son autonomie épistémologique pour devenir une discipline ancillaire, un « supplément d'âme » culturel.\cite{I-3-2-ref7} Cette approche nie la spécificité du latin comme langue système ayant sa propre logique.\item \textbf{La perte de la cohérence interne et la fragmentation des savoirs :} L'approche par projet (EPI) fragmente l'apprentissage. Au lieu d'une progression linguistique linéaire et cumulative — absolument nécessaire pour une langue flexionnelle où la maîtrise des cas (nominatif, accusatif, etc.) conditionne la suite — l'élève est confronté à des îlots de savoirs disparates. On étudie le vocabulaire du théâtre une semaine, celui de la citoyenneté le mois suivant, sans nécessairement construire la charpente morphologique (déclinaisons, conjugaisons) qui relie le tout. Cette fragmentation empêche la construction d'une compétence linguistique durable.\end{itemize}

\subsection{Critique des manuels de la réforme (2016-2017) : le syndrome du « magazine »}

Les manuels scolaires publiés dans la précipitation de la réforme (chez Hatier, Magnard, Nathan,
Belin) témoignent matériellement de cette dilution.\cite{I-3-2-ref18} L'analyse de ces ouvrages
révèle comment les éditeurs ont tenté de s'adapter à la quadrature du cercle imposée par les
programmes.

\begin{itemize}\item \textbf{Approche « magazine » et ludification :} Les manuels comme \textit{LCA Latin Cycle 4} (Hatier) ou \textit{Via Latina} multiplient les documents iconographiques en pleine page, les encadrés « Le saviez-vous? », les reconstitutions 3D et les passerelles avec la culture populaire contemporaine (Harry Potter, Percy Jackson, Hunger Games).\cite{I-3-2-ref21} Si cela peut séduire au premier abord et « accrocher » l'élève, cela réduit considérablement la place physique accordée aux textes latins authentiques et, surtout, aux tableaux de grammaire et d'exercices structuraux. Le manuel devient un livre d'images culturelles plus qu'un outil d'apprentissage de la langue.\item \textbf{Simplification et réécriture des textes :} Les textes proposés sont souvent des versions simplifiées, réécrites (« latin forcé ») ou massivement tronquées, accompagnées de traductions juxtalinéaires ou de traductions à trous. L'élève est placé en position de spectateur du texte plutôt qu'en acteur de sa traduction. On observe parfois des juxtapositions douteuses, mettant sur le même plan un texte de Voltaire et un extrait de \textit{Divergente} 22, brouillant les repères hiérarchiques littéraires.\item \textbf{Disparité des niveaux et absence de progression :} Conçus pour couvrir l'ensemble du cycle 4 (5ème, 4ème, 3ème) en un seul volume (pour des raisons économiques liées à la réforme), ces manuels peinent à offrir une progression claire et graduelle. Un élève de 3ème peut se retrouver à faire des exercices de niveau 5ème, faute de temps pour avoir acquis les bases intermédiaires, ou inversement être confronté à des textes trop complexes sans étayage suffisant.\cite{I-3-2-ref18} La grammaire est souvent reléguée en fin d'ouvrage sous forme d'annexes, comme un outil de consultation optionnel plutôt que comme le cœur de l'apprentissage.\end{itemize}C'est au cœur de la salle de classe que se noue le drame pédagogique de la réforme de 2016. Alors
que les programmes officiels et la réduction des horaires incitaient théoriquement à une approche «
culturelle » légère, la réalité des pratiques enseignantes a vu la persistance, voire le
renforcement paradoxal, de la méthode la plus traditionnelle et la plus coûteuse cognitivement : la
méthode grammaire-traduction.



\subsection{Le paradoxe temporel : enseigner comme au xixe siècle avec les horaires du xxie}

La méthode grammaire-traduction, héritée du XIXe siècle et de l'enseignement des humanités
classiques, repose sur un principe déductif : l'apprentissage explicite et par cœur des règles
(déclinaisons, conjugaisons, règles de syntaxe) et leur application systématique par la traduction
(thème et version).\cite{I-3-2-ref23} Cette méthode exige du temps, de la répétition, et une
exposition fréquente pour permettre l'assimilation.

Or, avec 1 heure en 5ème ou 2 heures en 4ème, cette méthode devient factuellement impraticable.
Pourtant, elle perdure massivement. Pourquoi cette inertie?

\begin{itemize}\item \textbf{L'habitus professionnel et la formation des enseignants :} La majorité des professeurs de Lettres Classiques en poste en 2016 ont été formés à l'université par cette méthode et ont passé les concours (CAPES, Agrégation) sur des épreuves de version et de thème académiques.\cite{I-3-2-ref24} En l'absence d'une formation continue massive et efficace aux méthodes actives (oralisation, méthode naturelle), ils reproduisent logiquement le modèle qu'ils maîtrisent et qu'ils identifient à la « rigueur » de la discipline. Ils tentent ainsi de faire entrer « au chausse-pied » le programme de grammaire traditionnel dans des horaires réduits, créant une densité de cours insupportable.\item \textbf{Sécurisation par la grammaire face à l'hétérogénéité :} Face à des classes de plus en plus hétérogènes et à des élèves parfois en difficulté de lecture ou de concentration, le cours de grammaire (remplir des tableaux, apprendre des terminaisons, faire des exercices à trous) offre un cadre rassurant, structuré et facilement évaluable pour l'enseignant. À l'inverse, les approches culturelles, de projet ou orales peuvent sembler plus floues à gérer, plus bruyantes, et moins légitimes aux yeux d'une profession attachée à la transmission d'un savoir technique précis.\cite{I-3-2-ref18}\end{itemize}

\subsection{La surcharge cognitive : analyse à la lumière des théories de sweller et tricot}

L'échec de cette rémanence méthodologique ne relève pas seulement d'un conservatisme pédagogique,
mais peut être analysé objectivement à travers le prisme de la psychologie cognitive, et notamment
la théorie de la charge cognitive développée par John Sweller et relayée en France par André
Tricot.\cite{I-3-2-ref26}



\subsection{Le coût cognitif intrinsèque du décodage latin}

Le latin est une langue à forte charge cognitive intrinsèque pour un apprenant francophone.
Contrairement à l'anglais ou à l'espagnol, qui partagent un ordre des mots (Sujet-Verbe-Objet)
relativement similaire au français, le latin est une langue flexionnelle où la fonction des mots est
déterminée par leur désinence (la fin du mot) et non par leur place. Le système casuel
(déclinaisons) oblige l'élève à traiter simultanément pour chaque mot :

1. La forme lexicale (le sens du mot).

2. La forme morphologique (la terminaison).

3. La fonction syntaxique associée à cette terminaison (Sujet, COD, Complément du nom, etc.).

Selon la théorie de la charge cognitive :

\begin{itemize}\item \textbf{Mémoire de travail limitée :} La mémoire de travail humaine ne peut traiter qu'un nombre très limité d'éléments nouveaux en interaction (généralement entre 4 et 7).\item \textbf{Nécessité de l'automatisation :} Pour réussir une tâche complexe comme la traduction ou la compréhension d'un texte, les tâches de bas niveau (reconnaissance des cas, identification du vocabulaire) doivent être automatisées (stockées en mémoire à long terme) pour libérer de la ressource mentale dans la mémoire de travail.\cite{I-3-2-ref31} Si l'élève doit réfléchir consciemment à chaque terminaison (« est-ce un accusatif ou un ablatif? »), sa mémoire de travail sature instantanément, et il perd le fil du sens global de la phrase.\end{itemize}

\subsection{La « double peine » cognitive de la réforme}

La réforme de 2016 a créé une situation de « double peine » cognitive pour les élèves :

1. \textbf{L'impossibilité de l'automatisation :} Avec 1h ou 2h par semaine, l'automatisation des
mécanismes de bas niveau est impossible. La courbe de l'oubli fait son œuvre d'une semaine sur
l'autre. À chaque cours, l'élève doit réapprendre ou « redécouvrir » les déclinaisons. Les
connaissances ne passent jamais en mémoire à long terme procédurale. Le coût cognitif reste donc
maximal à chaque séance.

2. \textbf{L'injonction contradictoire (Top-down vs Bottom-up) :} D'un côté, les programmes
demandent aux élèves de gérer des tâches complexes de « repérage d'indices » culturels et
linguistiques dans des textes authentiques (approche \textit{top-down} ou descendante). De l'autre,
faute d'automatisation des bases grammaticales (approche \textit{bottom-up} ou ascendante), chaque
mot latin reste une énigme indéchiffrable nécessitant un traitement conscient et coûteux.

Le résultat est une surcharge cognitive permanente. L'élève ne lit pas le latin ; il le déchiffre
péniblement, mot à mot, comme un code secret, en se référant constamment à des tableaux de grammaire
annexes qu'il ne maîtrise pas. Cette méthode, appliquée sur des horaires homéopathiques, conduit
inévitablement au découragement, au sentiment d'incompétence et à l'échec pour tous les élèves qui
ne bénéficient pas d'un soutien extérieur familial (capital culturel) pour compenser les manques de
l'école. La réforme a ainsi renforcé la reproduction sociale qu'elle prétendait combattre : seuls
ceux qui ont les codes culturels ou le soutien privé peuvent réussir dans ce système dysfonctionnel.



\subsection{L'échec de la transition vers les méthodes actives : l'exemple du projet ellass}

Face à ce constat d'échec de la méthode traditionnelle, des alternatives didactiques existent
pourtant, prônant une pédagogie active inspirée des méthodes d'enseignement des langues vivantes
(oralisation, méthode naturelle/Ørberg, méthode audio-orale). Le projet de recherche action LéA «
ELLASS » (Enseigner les Langues Anciennes dans le Secondaire et le Supérieur) incarne cette
tentative de renouveau didactique.\cite{I-3-2-ref32}



\subsection{Les promesses de l'oralisation et de la paraphrase progressive}

Le projet ELLASS, documenté par l'Institut Français de l'Éducation (IFÉ), propose des pistes
audacieuses pour contourner l'obstacle de la grammaire-traduction :

\begin{itemize}\item \textbf{L'oralisation :} Traiter le latin et le grec comme des langues de communication (approche communicative). Faire parler les élèves, utiliser des dialogues, pour ancrer les structures grammaticales par l'oreille et l'usage plutôt que par l'analyse abstraite.\item \textbf{La paraphrase progressive :} Utiliser des textes simplifiés et hiérarchisés (paraphrases) pour permettre aux élèves d'accéder au sens du texte littéraire progressivement, sans passer par la traduction mot à mot immédiate.\cite{I-3-2-ref32} Cela réduit la charge cognitive initiale et permet une entrée plus fluide dans la langue.\item \textbf{L'approche inductive :} Partir de l'exemple et de l'usage pour déduire la règle, plutôt que l'inverse.\end{itemize}

\subsection{L'incompatibilité structurelle avec les horaires de 2016}

Cependant, ces méthodes actives se heurtent à un mur de réalité : elles nécessitent une fréquence
d'exposition élevée pour être efficaces. L'immersion linguistique, l'oralisation, l'imprégnation par
la méthode naturelle (type Ørberg) requièrent idéalement une pratique quotidienne ou
plurigebdomadaire (3 à 4 heures minimum).

La structure horaire de 2016 (1h à 2h) rend l'oralisation anecdotique, voire folklorique. On ne peut
pas « apprendre à parler latin » ou s'imprégner d'une langue avec une séance de 50 minutes par
semaine. Les enseignants innovants qui tentent ces méthodes se retrouvent souvent bloqués par le
manque de temps de pratique, tandis que la majorité se replie sur la grammaire-traduction par
défaut, jugée plus « rentable » à court terme pour préparer aux évaluations
écrites.\cite{I-3-2-ref33} La réforme a ainsi structurellement empêché la modernisation pédagogique
qu'elle appelait de ses vœux.

La réforme de 2016 a initié une dynamique délétère qui s'est poursuivie et aggravée sous les
ministères suivants, malgré des ajustements de façade. Loin de sauver les LCA, elle a organisé leur
marginalisation progressive, les rendant vulnérables aux coups de boutoir ultérieurs.



\subsection{L'illusion du redressement (2017-2022) : le temps de la survie précaire}

L'arrivée de Jean-Michel Blanquer au ministère en 2017 a semblé apporter un répit politique. Le
rétablissement symbolique de l'option LCA (distincte des EPI, qui ne sont plus obligatoires pour le
latin) et la réintroduction de la possibilité théorique de remonter à 3 heures en 3ème ont été
salués par certains comme des « signaux positifs ».\cite{I-3-2-ref1}

Cependant, le mécanisme de financement (la marge d'autonomie) n'a pas été modifié.

\begin{itemize}\item \textbf{La concurrence accrue :} Les établissements ont dû financer de nouveaux dispositifs prioritaires (Devoirs Faits, dédoublements des classes de CP/CE1 qui ont impacté les moyens globaux, puis les groupes de besoins) sur la même enveloppe de marge. Les LCA sont restées la variable d'ajustement idéale.\item \textbf{L'effet d'inertie :} De nombreux collèges, ayant basculé sur le modèle 1h/2h/2h en 2016, l'ont maintenu par inertie administrative et confort de gestion, entérinant la baisse de niveau et d'exigence.\cite{I-3-2-ref1} L'exception est devenue la norme.\end{itemize}

\subsection{Le coup de grâce : le « choc des savoirs » (2024)}

La réforme dite du « Choc des Savoirs », initiée par Gabriel Attal et mise en œuvre pour la rentrée
2024, constitue l'aboutissement logique et peut-être fatal de la marginalisation initiée en
2016.\cite{I-3-2-ref2}



\subsection{Les groupes de niveaux comme prédateurs absolus de moyens}

La mise en place des « groupes de niveaux » (ou groupes de besoins) en français et en mathématiques
nécessite un redéploiement massif et inédit des heures de la marge d'autonomie pour créer les
barrettes de cours supplémentaires.

\begin{itemize}\item \textbf{Siphonnage de la marge :} Pour créer des groupes réduits en français/maths (priorité nationale absolue), les chefs d'établissement sont contraints mathématiquement de récupérer toutes les heures disponibles. Les options facultatives comme les LCA sont les premières victimes. Le SNES et la CNARELA rapportent pour la rentrée 2024 des suppressions massives d'heures de LCA, des refus d'ouverture en 5ème, ou des regroupements de niveaux pédagogiquement aberrants (élèves de 5ème, 4ème et 3ème mélangés dans la même heure de cours) pour sauver l'existence administrative de l'option.\cite{I-3-2-ref2}\item \textbf{L'exception SEGPA :} Même les dispositifs pour les élèves les plus fragiles sont touchés, créant une tension sur l'ensemble des moyens.\cite{I-3-2-ref39}\end{itemize}

\subsection{Une discipline en voie d'extinction rapide?}

Les rapports syndicaux (SNES, FO) et associatifs (Arrête Ton Char) décrivent une situation
catastrophique pour 2024-2025. La discipline est menacée de disparition pure et simple dans de
nombreux collèges publics, faute de moyens. Les professeurs de Lettres Classiques se retrouvent
contraints de partager leur service sur trois ou quatre établissements (TZR) pour compléter leurs
heures, ou de n'enseigner que du français moderne, perdant leur spécificité.\cite{I-3-2-ref2}

La « fausse solution » de 2016 a ainsi préparé le terrain à une extinction de fait : en rendant la
discipline optionnelle, mal financée, pédagogiquement floue et dépendante des arbitrages locaux,
elle l'a rendue indéfendable face aux nouvelles priorités ministérielles basées sur les « savoirs
fondamentaux » (lire, écrire, compter).



\subsection{Bilan qualitatif final : quel latin pour quels élèves en 2026?}

Au terme de cette analyse, une question demeure : qu'apprennent réellement les élèves qui suivent
encore ces cours « homéopathiques » rescapés des réformes?

\begin{itemize}\item \textbf{Une culture de la reconnaissance (passive) :} Ils apprennent à reconnaître des mythes, des racines étymologiques, des monuments. C'est une culture de l'honnête homme, utile pour la culture générale, mais souvent superficielle et déconnectée du texte source.\item \textbf{Une ignorance linguistique (structurelle) :} La majorité des élèves sortent du collège sans savoir lire un texte latin simple sans dictionnaire et traduction juxtalinéaire. La compétence de traduction, qui structure la pensée logique et l'analyse syntaxique (le fameux « esprit critique » et la rigueur vantés par les défenseurs des humanités), est la grande perdante de la décennie.\item \textbf{Un élitisme résiduel renforcé :} Paradoxalement, la réforme a renforcé l'élitisme. Seuls les élèves les plus favorisés, ou ceux scolarisés dans des établissements privés ou des collèges de centre-ville privilégiés ayant conservé des horaires décents sur fonds propres (parfois financés par les familles ou des fondations), accèdent encore à la véritable maîtrise linguistique. La réforme de 2016, censée démocratiser, a recréé une distinction brutale entre un « latin culturel » (EPI/Option dégradée) pour le peuple et un « latin linguistique » (Option renforcée/Privé) pour l'élite.\end{itemize}L'analyse approfondie de la situation des Langues et Cultures de l'Antiquité depuis la réforme de
2016 permet de valider sans ambiguïté la thèse de la « fausse solution ». La réforme n'a pas sauvé
les langues anciennes ; elle a organisé leur survie artificielle sous une forme dégradée, en
organisant leur dilution.

\textbf{La dilution dans la culture} a vidé la discipline de sa substance rigoureuse, remplaçant
l'exercice formateur de la traduction par une approche documentaire souvent passive et une
intertextualité de surface. \textbf{La rémanence de la méthode grammaire-traduction}, faute de
temps, de moyens et de formation pour implémenter des alternatives pédagogiques viables (méthodes
actives, oralisation), a enfermé les élèves et les enseignants dans une impasse cognitive : essayer
de faire entrer une cathédrale grammaticale dans un studio d'une heure par semaine. Cette surcharge
cognitive a généré de l'échec et du découragement.

Enfin, les \textbf{horaires homéopathiques} et le mécanisme pervers du financement sur la marge
d'autonomie ont transformé le latin et le grec en variables d'ajustement comptable, les livrant sans
défense aux appétits des réformes ultérieures comme le « Choc des savoirs » de 2024. En 2026,
l'enseignement des langues anciennes au collège public français apparaît ainsi comme un astre mort :
il brille encore théoriquement dans les programmes officiels, mais sa réalité matérielle et
pédagogique s'est éteinte pour le plus grand nombre, victime d'une réforme qui a tragiquement
confondu l'affichage politique de la démocratisation avec la réalité exigeante de la transmission
des savoirs.



\subsection{Tableaux de synthèse des données}



\subsection{Tableau 1 : évolution comparée des horaires officiels de latin (cycle 4\)}

\begin{table}[h!]\small\centering\begin{tabular}{| l | l | l | l | l |}\hline\textbf{Niveau} & \textbf{Avant 2016 (Horaires fléchés)} & \textbf{Réforme 2016 (Arrêté 2015)} & \textbf{Ajustement 2017 (\og Dans la limite de\fg{})} & \textbf{Réalité constatée 2024 (SNES/CNARELA)} \\ \hline\textbf{5ème} & 2h (fréquent) & 1h & 1h & 1h (souvent globalisée ou supprimée) \\ \hline\textbf{4ème} & 3h & 2h & 3h (max) & 2h majoritaire \\ \hline\textbf{3ème} & 3h & 2h & 3h (max) & 2h ou regroupement avec 4ème \\ \hline\textbf{Total Cycle} & \textbf{8h} & \textbf{5h} & \textbf{7h (théorique)} & \textbf{\~4h à 5h réelles} \\ \hlineSources croisées : 1 &  &  &  &  \\ \hline\end{tabular}\end{table}

\subsection{Tableau 2 : comparaison des approches didactiques et impact cognitif}

\begin{table}[h!]\small\centering\begin{tabular}{| p{2cm} | p{2.2cm} | p{2.2cm} | p{2.2cm} |}\hline\textbf{Critère} & \textbf{Méthode Traditionnelle (Grammaire-Traduction)} & \textbf{Approche Réforme 2016 (\og Culturelle\fg{})} & \textbf{Approche Active (Projet ELLASS)} \\ \hline\textbf{Objectif} & Maîtrise du code, Version, Thème & Repérage d'indices, Civilisation, Étymologie & Communication, Lecture cursive, Oralisation \\ \hline\textbf{Support} & Textes d'auteurs, Tableaux grammaticaux & Documents composites, Textes traduits, Images & Textes paraphrasés, Dialogues, Audio \\ \hline\textbf{Charge Cognitive} & \textbf{Très élevée} (Décodage explicite constant) & \textbf{Faible} (Approche globale/devinette) & \textbf{Progressive} (Inductive, étayée) \\ \hline\textbf{Compatibilité Horaires 2016} & \textbf{Impossible} (Manque de temps d'automatisation) & \textbf{Possible} (Mais faible gain linguistique) & \textbf{Difficile} (Manque de fréquence d'exposition) \\ \hlineSources croisées : 13 &  &  &  \\ \hline\end{tabular}\end{table}